\scanpage{5}
\begin{center}Sur les groupes de Galois attachés aux groupes $p$-divisibles\hfill119\end{center}

naturelle de groupe de Lie $p$-adique. Soit $G$ un sous-groupe compact de $\operatorname{Aut}(V)$. On sait (voir par exemple [7], p. 5.42) que $G$ est un \textit{sous-groupe de Lie} de $\operatorname{Aut}(V)$. Soit $\mathfrak g$ son algèbre de Lie; l'exponentielle définit un isomorphisme d'un voisinage ouvert de 0 dans $\mathfrak g$ (muni de la loi de groupe fournie par la formule de Hausdorff) sur un sous-groupe ouvert de $G$.

\textbf{Proposition 1.} -- \textit{Les conditions suivantes sont équivalentes:}

(a) \textit{$V$ est un $G$-module semi-simple.}

(b) \textit{$V$ est un $\mathfrak g$-module semi-simple.}

(c) \textit{$\mathfrak g$ est une algèbre de Lie réductive} (i.e. produit d'une algèbre abélienne $\mathfrak c$ par une algèbre semi-simple $\mathfrak s$) \textit{et $V$ est un $\mathfrak c$-module semi-simple.}

Soit $G_1$ un sous-groupe ouvert distingué de $G$ assez petit pour être contenu dans l'image de l'exponentielle. Tout sous-espace de $V$ stable par $\mathfrak g$ l'est aussi par $G_1$, et inversement. Cela montre l'équivalence de (b) et de la condition suivante:

(a$_1$) $V$ est un $G_1$-module semi-simple.

Mais $G/G_1$ est fini. On en déduit facilement (cf. \textsc{Chevalley} [1], tome III, p. 82, prop. 1) l'équivalence de (a) et (a$_1$). D'où $(\mathrm a)\Leftrightarrow(\mathrm b)$. L'équivalence $(\mathrm b)\Leftrightarrow(\mathrm c)$ est bien connue (\textsc{Bourbaki}, \textit{Gr. et Alg. de Lie}, chap. I, \S\ 6, n\textsuperscript{o} 5, th. 4).

Nous supposerons à partir de maintenant que les conditions de la prop. 1 sont vérifiées. Nous allons associer à $G$ un certain groupe algébrique réductif $G_{\mathrm{alg}}$, de la manière suivante:

Soit d'abord $\mathbf{GL}_V$ le $\mathbf Q_p$-groupe algébrique des automorphismes de $V$. Cela signifie, par définition, que, si $k$ est une $\mathbf Q_p$-algèbre commutative, le groupe $\mathbf{GL}_V(k)$ des points de $\mathbf{GL}_V$ à valeurs dans $k$ est égal à $\operatorname{Aut}(V\otimes k)$. En particulier, le groupe $\mathbf{GL}_V(\mathbf Q_p)$ s'identifie au groupe $\operatorname{Aut}(V)$ considéré plus haut.

\textbf{Définition.} -- \textit{On appelle enveloppe algébrique de $G$, et on note $G_{\mathrm{alg}}$, le plus petit sous-groupe algébrique de $\mathbf{GL}_V$ dont le groupe des points contienne $G$.}

L'existence et l'unicité de $G_{\mathrm{alg}}$ sont immédiates (et ne nécessitent aucune hypothèse sur $G$). Si $A$ est l'algèbre affine de $\mathbf{GL}_V$, l'idéal définissant $G_{\mathrm{alg}}$ est l'ensemble des $f\in A$ tels que $f(g)=0$ pour tout $g\in G$.

Utilisons maintenant l'hypothèse faite sur $G$. Décomposons l'algèbre de Lie $\mathfrak g$ en
\[
\mathfrak g=\mathfrak c\times\mathfrak s,
\]
où $\mathfrak c$ est le centre de $\mathfrak g$, et $\mathfrak s=[\mathfrak g,\mathfrak g]$ est semi-simple. On sait (\textsc{Chevalley}
\margin{4}
